% Earth-orbit and lunar-orbit families at the same scale (1 cm = 22 000 km).
% Ellipses are two-body schematics; the NRHO and the L2 halo are three-body orbits whose shapes are indicated, not propagated.
\begin{tikzpicture}[x=1cm,y=1cm,font=\scriptsize]
% ===== Earth panel =====
\begin{scope}[shift={(4.6,0)}]
\node[font=\footnotesize\bfseries,anchor=south] at (0,4.55) {(a) Earth orbits};
\fill[blue!55!cyan!70!black] (0,0) circle (0.290);
\begin{scope}\clip (0,0) circle (0.290);
\fill[green!45!black!80] (-0.20,0.05) .. controls (-0.12,0.20) and (0.02,0.16) .. (0.04,0.04) .. controls (0.06,-0.08) and (-0.05,-0.20) .. (-0.12,-0.12) .. controls (-0.18,-0.06) and (-0.24,-0.03) .. cycle;
\fill[green!40!brown!80] (0.10,0.12) .. controls (0.18,0.20) and (0.26,0.10) .. (0.22,0.0) .. controls (0.18,-0.10) and (0.12,-0.02) .. cycle;
\fill[white!90!blue] (-0.12,0.25) rectangle (0.12,0.30);
\fill[white,opacity=0.35] (-0.09,0.10) circle (0.10);
\end{scope}
\draw[black!60] (0,0) circle (0.290);
\draw[black!70,thick] (0,0) circle (0.315); \draw[black!40,thin] (-0.297,-0.105) -- (-2.25,-0.75); \node[anchor=east,text=black!70] at (-2.3,-0.75) {LEO, \SI{600}{km}};
\draw[black!70] (0,0) circle (1.207); \draw[black!40,thin] (-0.925,0.776) -- (-2.25,1.45); \node[anchor=east,text=black!70] at (-2.3,1.45) {MEO (GNSS), \SI{20200}{km}};
\draw[black!70,thick] (0,0) circle (1.917); \node[anchor=south,text=black!70] at (0,1.967) {GEO, \SI{35786}{km}};
\draw[black!50,dashed,rotate=63] (0.895,0) ellipse (1.209 and 0.813);
\node[anchor=west,text=black!50,align=left] at (1.15,2.05) {Molniya,\\$500\times\SI{39800}{km}$};
\fill[black!70] (0,0.315) circle (0.03); \fill[black!70] (0,1.917) circle (0.04); \fill[black!70] (-0.925,0.776) circle (0.03);
\draw[-{Stealth[length=4pt]},black!60] (3.3,-4.45) -- (4.3,-4.45); \node[anchor=east,text=black!60] at (3.3,-4.45) {to the Moon, \SI{384400}{km} (17 cm at this scale)};
\end{scope}
% ===== Moon panel =====
\begin{scope}[shift={(13.2,0)}]
\node[font=\footnotesize\bfseries,anchor=south] at (0,4.55) {(b) Lunar orbits, same scale};
\draw[black!35,densely dotted] (0,0) circle (2.727); \node[anchor=south,text=black!45] at (0,2.76) {Hill sphere, ${\approx}\SI{60000}{km}$};
\draw[red!60,densely dotted] (0,0) circle (4.018); \node[anchor=south,text=red!70,align=center] at (0,4.05) {synchronous radius \SI{88400}{km}, outside the Hill sphere};
\fill[black!22] (0,0) circle (0.079);
\fill[black!38] (-0.025,0.02) circle (0.022); \fill[black!38] (0.03,-0.025) circle (0.015);
\draw[black!60] (0,0) circle (0.079);
\draw[black!70,thick] (0,0) circle (0.083); \node[anchor=south west,text=black!70] at (0.12,0.08) {LLO, \SI{100}{km}};
% ELFO: apolune south (down)
\draw[green!45!black,thick] (0,-0.356) ellipse (0.371 and 0.514);
\draw[green!45!black,thin] (-0.33,-0.20) -- (-1.25,0.75); \node[anchor=south east,text=green!45!black,align=right] at (-1.25,0.75) {ELFO (LCRNS)\\$1750\times\SI{17400}{km}$, \SI{30}{h}};
% NRHO schematic: apolune south
\draw[brown!70!black,thick,dash pattern=on 4pt off 1.5pt] (0,-1.539) ellipse (0.703 and 1.692);
\draw[brown!70!black,thin] (0.62,-2.25) -- (2.85,-3.0); \node[anchor=west,text=brown!70!black] at (2.9,-3.0) {NRHO (Gateway), $1600\times\SI{70000}{km}$, \SI{6.6}{d}};
% L2 point and halo, anti-Earth direction = +x
\fill[black!60] (2.795,0) circle (0.04); \node[anchor=south,text=black!60] at (2.795,0.6) {L2 halo (Queqiao)};
\draw[blue!50!black,thick] (2.795,0) ellipse (0.3 and 0.6);
\draw[-{Stealth[length=4pt]},black!60] (-3.0,-4.45) -- (-4.0,-4.45); \node[anchor=west,text=black!60,align=left] at (-3.0,-4.45) {to Earth, \SI{384400}{km}: from the near side,\\Earth is nearly fixed in the sky};
\end{scope}
\end{tikzpicture}
